%%%%%%%%%%%%%%%%%%%%%%%%%%%%%%%%%
%
%       EXERCÍCIO
%
%%%%%%%%%%%%%%%%%%%%%%%%%%%%%%%%%

\ifdefstring{\atividade}{prova}{%
    \ifdefstring{\modo}{objetivo}{%
        \renewcommand{\valorquestao}{\ValorQObj\ ponto}
    }{%
        \renewcommand{\valorquestao}{\ValorQDisc\ pontos}
    }%
}%

\begin{exercicioBanco}[\valorquestao]
Dois reservatórios de água, inicialmente com o mesmo volume de $20$ litros, começam a ser abastecidos. Os volumes de água, $V_{1}$ e $V_{2}$, em função do tempo $t$, em minutos, estão representados no plano cartesiano pelas retas com interseção no ponto $A = (0,20)$ e que passam, respectivamente, pelos pontos $B = (10, 50)$ e $C = (30, 35)$.

\begin{center}
    \begin{tikzpicture}[scale=0.5]
        % Eixos
        \draw[->] (0,0) -- (6,0) node[below right] {\small{$t$ (min)}};
        \draw[->] (0,0) -- (0,6) node[above left] {\small{$V$ (L)}};
        
        % Funções
        \draw[domain=0:1.5, smooth, variable=\x, blue, thick] plot ({\x}, {3*\x+2}) node[below right] {\small $V_{1}$};
        \draw[domain=0:5, smooth, variable=\x, red, thick] plot ({\x}, {0.5*\x+2}) node[below right] {\small $V_{2}$};
        
        % Pontos
        \filldraw (0,2) circle (2pt) node[left] {\small $A$};
        \filldraw (1,5) circle (2pt) node[below right] {\small $B$};
        \filldraw (3,3.5) circle (2pt) node[below right] {\small $C$};
    \end{tikzpicture}
\end{center}

Sabe-se que, em determinado intervalo de tempo, o volume $V_{2}$ aumentou $5$ litros.

Nesse mesmo intervalo de tempo, determine quanto aumentou o volume $V_{1}$, em litros.


% Define as alternativas
\newcommand{\alternativas}{%
    \begin{center}
        \begin{tabularx}{\textwidth}{XXXXX}
            (a) \(10\). &
            (b) \(15\). &
            (c) \(20\). &
            (d) \(25\). &
            (e) \(30\).
        \end{tabularx}
    \end{center}
}

% Resposta correta
\newcommand{\resposta}{E}

% Lógica condicional para exibição
\ifdefstring{\atividade}{lista}{%
    \alternativas
    \vspace{0.5em}
    
    \noindent\textbf{Resposta:} letra \textbf{\resposta}.
}{%
    \ifdefstring{\modo}{objetivo}{%
        \alternativas
    }{%
        % modo = discursiva → não mostra alternativas
    }
}
\end{exercicioBanco}

